% Generated draft: requires mathematical and source review.
\documentclass[11pt]{article}
\usepackage[margin=25mm]{geometry}
\usepackage{fontspec}
\setmainfont{Latin Modern Roman}
\usepackage{amsmath,amssymb,mathtools,unicode-math}
\setmathfont{Latin Modern Math}
\usepackage{enumitem,needspace,xurl,hyperref}
\setlength{\parindent}{0pt}
\setlength{\parskip}{5pt}
\setlength{\emergencystretch}{3em}
\newcommand{\N}{\mathbb{N}}
\newcommand{\Z}{\mathbb{Z}}
\newcommand{\Q}{\mathbb{Q}}
\newcommand{\R}{\mathbb{R}}
\newcommand{\C}{\mathbb{C}}
\newcommand{\F}{\mathbb{F}}
\newcommand{\A}{\mathbb{A}}
\newcommand{\PP}{\mathbb{P}}
\newcommand{\E}{\mathbb{E}}
\newcommand{\eps}{\varepsilon}
\newcommand{\dd}{\,\mathrm{d}}
\begin{document}
\section*{UnsolvedMath ID: 701279}
\textbf{Source problem code:} catalogue entry
\subsection{Definitions and mathematical statement}
Let \(G\) be a group with identity \(1_G\). A subset \(X\subseteq G\) is a generating set if \(\langle X\rangle=G\). For such \(X\), define
\[
\ell_X(g)=\min\bigl\{n\in\{0,1,2,\ldots\}:g=x_1^{\varepsilon_1}\cdots x_n^{\varepsilon_n},\ x_i\in X,\ \varepsilon_i\in\{-1,1\}\bigr\},
\]
where the empty product has length zero and equals \(1_G\). Thus group words may use both generators and their inverses. The width of \(G\) with respect to \(X\) is
\[
\operatorname{width}_X(G)=\sup_{g\in G}\ell_X(g)\in\{0,1,2,\ldots\}\cup\{\infty\}.
\]
Following the supplied source, \(G\) has \emph{finite width} if every generating subset \(X\subseteq G\), including infinite generating subsets, gives finite width. Equivalently,
\[
\forall X\subseteq G\quad\bigl(\langle X\rangle=G\ \Longrightarrow\ \exists N_X\in\{0,1,2,\ldots\}\ \forall g\in G:\ell_X(g)\leq N_X\bigr).
\]
The bound \(N_X\) may depend on \(X\).

Does there exist a group \(G\) with \(|G|=\aleph_0\) that has finite width?

The supplied source reports that all known infinite groups of finite width are uncountable, mentioning infinite permutation groups, infinite-dimensional general linear groups, and some other groups. This is preserved as a source-specific literature qualification, without asserting that every group in either named class has finite width or that the assessment remains current.
\subsection{Short English statement}
Does there exist a countably infinite group such that, for each generating set, every element can be expressed using a bounded number of generators and their inverses? The bound may depend on the generating set. The supplied registry labels this problem open; that label is not independently verified here.
\subsection{Sources}
\begin{itemize}
\item Supplied local UnsolvedMath registry record, permanent ID \(701279\), exact source problem code \texttt{catalogue entry}, titled \emph{Kourovka Notebook Problem 16.88}. The supplied status is \texttt{open}.
\item \emph{The Kourovka Notebook (Unsolved Problems in Group Theory)}, problem 16.88, 16th issue (2006); the supplied background also identifies the 21st issue edition (2026), arXiv:1401.0300v48. The collection description names E. I. Khukhro and V. D. Mazurov as editors.
\item The registry attributes the proposal to V. Tolstykh, with proposed year 2006. The statement itself begins with the attribution ``G. M. Bergman.'' These attribution roles are retained separately.
\item G. M. Bergman, \emph{Bull. London Math. Soc.} \textbf{38} (2006), 429--440, and references therein, as cited in the supplied statement for the background on infinite groups of finite width. No article title or external link is supplied for this reference.
\end{itemize}
\end{document}
